\section{Базис и размерность линейного пространства. Основная лемма о линейной зависимости. Теорема о связи базиса и размерности}

\begin{definition}
Линейное пространство $V$ называется \textbf{конечномерным} размерности $n$ ($\dim V = n$), если:
\begin{enumerate}
    \item В $V$ существует система из $n$ линейно независимых векторов.
    \item Любая система из $n+1$ вектора в $V$ линейно зависима.
\end{enumerate}
Если такого числа $n$ не существует, пространство называется бесконечномерным ($\dim V = \infty$).
\end{definition}

\begin{definition}
Упорядоченная система векторов $E = \{e_1, e_2, \dots, e_k\}$ называется \textbf{базисом} линейного пространства $V$, если:
\begin{enumerate}
    \item Векторы $e_1, \dots, e_k$ линейно независимы.
    \item Любой вектор $x \in V$ может быть представлен в виде их линейной комбинации: $x = \xi_1 e_1 + \xi_2 e_2 + \dots + \xi_k e_k$, где $\xi_i \in F$ --- координаты вектора $x$ в базисе $E$.
\end{enumerate}
\end{definition}

\begin{lemma}[Основная лемма о линейной зависимости]
Пусть $S_1 = \{a_1, a_2, \dots, a_k\}$ и $S_2 = \{b_1, b_2, \dots, b_m\}$ --- две системы векторов из $V$. Если система $S_1$ линейно независима и каждый вектор из $S_1$ линейно выражается через векторы системы $S_2$, то количество векторов в первой системе не превосходит количества векторов во второй, то есть $k \le m$.
\end{lemma}
\begin{remark}
Доказательство данной леммы приведено в билете №7.
\end{remark}

\begin{theorem}[О связи базиса и размерности]
\begin{enumerate}
    \item Если $\dim V = n$, то любая система из $n$ линейно независимых векторов образует базис в $V$.
    \item Если в пространстве $V$ существует базис, состоящий из $n$ векторов, то $\dim V = n$.
\end{enumerate}
\end{theorem}

\begin{proof}
\begin{enumerate}
    \item Пусть $E = \{e_1, \dots, e_n\}$ --- система из $n$ линейно независимых векторов в $V$. Возьмем произвольный вектор $x \in V$. Рассмотрим систему $\{e_1, \dots, e_n, x\}$, состоящую из $n+1$ вектора. Поскольку $\dim V = n$, любая система из $n+1$ вектора линейно зависима. Значит, система $\{e_1, \dots, e_n, x\}$ линейно зависима. 
    Так как подсистема $\{e_1, \dots, e_n\}$ по условию линейно независима, то по теореме о добавлении вектора (билет 11) вектор $x$ линейно выражается через $e_1, \dots, e_n$. Следовательно, $E$ удовлетворяет обоим условиям определения базиса.
    
    \item Пусть $E = \{e_1, \dots, e_n\}$ --- базис $V$. По определению базиса, эта система из $n$ векторов линейно независима (условие 1 размерности выполнено). Пусть существует система из $n+1$ вектора $\{x_1, \dots, x_{n+1}\}$. Каждый из этих векторов линейно выражается через базис $E$ (из $n$ векторов). По основной лемме о линейной зависимости, так как $n+1 > n$, система $\{x_1, \dots, x_{n+1}\}$ обязана быть линейно зависимой (условие 2 размерности выполнено). Следовательно, $\dim V = n$.
\end{enumerate}
\hfill$\blacksquare$
\end{proof}

\begin{example}[Примеры базисов]
\begin{enumerate}
    \item В пространстве $\mathbb{R}^n$ стандартным базисом является система векторов $e_1 = (1, 0, \dots, 0)^T, \dots, e_n = (0, \dots, 0, 1)^T$.
    \item В пространстве матриц $M_{m \times n}$ базис образуют \textbf{матричные единицы} $E_{ij}$ (матрицы, у которых на пересечении $i$-й строки и $j$-го столбца стоит $1$, а все остальные элементы равны $0$). Таких матриц ровно $m \times n$, следовательно, $\dim M_{m \times n} = m \cdot n$.
    \[
    E_{ij} = \begin{pmatrix}
    0 & \dots & 0 & \dots & 0 \\
    \vdots & \ddots & \vdots & & \vdots \\
    0 & \dots & 1 & \dots & 0 \\
    \vdots & & \vdots & \ddots & \vdots \\
    0 & \dots & 0 & \dots & 0
    \end{pmatrix}
    \begin{matrix}
    \\ \\ \leftarrow i\text{-я строка} \\ \\
    \end{matrix}
    \]
\end{enumerate}
\end{example}
